
\documentclass{article}
\usepackage{colortbl}
\usepackage{makecell}
\usepackage{multirow}
\usepackage{supertabular}

\begin{document}

\newcounter{utterance}

\centering \large Interaction Transcript for game `hot\_air\_balloon', experiment `air\_balloon\_survival\_ru\_negotiation\_easy', episode 2 with deepseek{-}v4{-}pro.
\vspace{24pt}

{ \footnotesize  \setcounter{utterance}{1}
\setlength{\tabcolsep}{0pt}
\begin{supertabular}{c@{$\;$}|p{.15\linewidth}@{}p{.15\linewidth}p{.15\linewidth}p{.15\linewidth}p{.15\linewidth}p{.15\linewidth}}
    \# & $\;$A & \multicolumn{4}{c}{Game Master} & $\;\:$B\\
    \hline

    \theutterance \stepcounter{utterance}  
    & & \multicolumn{4}{p{0.6\linewidth}}{
        \cellcolor[rgb]{0.9,0.9,0.9}{
            \makecell[{{p{\linewidth}}}]{
                \texttt{\tiny{[P1$\langle$GM]}}
                \texttt{Вы участвуете в совместной переговорной игре.} \\
\\ 
\texttt{Вместе с другим участником вы должны договориться об одном наборе предметов, которые будут сохранены. У каждого из вас есть своё представление о том, насколько важен для вас каждый предмет (важность). Вы не знаете, как другой участник оценивает предметы. Кроме того, для каждого предмета указаны затраты, которые он требует.} \\
\texttt{Вы можете согласовать набор только в том случае, если суммарные затраты выбранных предметов не превышает общий предел:} \\
\\ 
\texttt{ПРЕДЕЛ = 3288} \\
\\ 
\texttt{Вот значения затрат для отдельных предметов:} \\
\\ 
\texttt{Затраты предмета = \{"C08": 329, "C55": 41, "A43": 417, "A67": 75, "B42": 390, "A62": 887, "B19": 808, "C43": 151, "C73": 849, "C10": 129, "B36": 350, "A18": 118, "C61": 630, "C67": 602, "C42": 801\}} \\
\\ 
\texttt{Вот ваше личное представление о важности каждого предмета:} \\
\\ 
\texttt{Значения важности предметов = \{"C08": 138, "C55": 583, "A43": 868, "A67": 822, "B42": 783, "A62": 65, "B19": 262, "C43": 121, "C73": 508, "C10": 780, "B36": 461, "A18": 484, "C61": 668, "C67": 389, "C42": 808\}} \\
\\ 
\texttt{Цель:} \\
\\ 
\texttt{Ваша цель — договориться об общем наборе предметов, который приносит вам как можно больше выгоды (то есть максимизирует суммарную важность для ВАС), оставаясь в ПРЕДЕЛАХ. Вы не обязаны делать ПРЕДЛОЖЕНИЕ в каждом сообщении — вы можете просто вести переговоры. Разрешены любые тактики!} \\
\\ 
\texttt{Протокол взаимодействия:} \\
\\ 
\texttt{В сообщении вы можете использовать только следующие структурированные форматы:} \\
\\ 
\texttt{ПРЕДЛОЖЕНИЕ: \{'А', 'Б', 'В', …\}} \\
\texttt{Предложите сохранить ровно эти предметы.} \\
\texttt{ОТКАЗ: \{'А', 'Б', 'В', …\}} \\
\texttt{Явно отклоните предложение оппонента.} \\
\texttt{АРГУМЕНТ: \{'...'\}} \\
\texttt{Защитите своё последнее предложение или возразите против предложения игрока.} \\
\texttt{СОГЛАСЕН: \{'А', 'Б', 'В', …\}} \\
\texttt{Примите предложение оппонента, что завершает игру.} \\
\\ 
\\ 
\texttt{Правила:} \\
\\ 
\texttt{Вы можете дать СОГЛАСЕН только на предложение, которое другая сторона зафиксировала через ПРЕДЛОЖЕНИЕ.} \\
\texttt{Вы можете дать ОТКАЗ только на предложение, которое другая сторона зафиксировала через ПРЕДЛОЖЕНИЕ.} \\
\texttt{Суммарное усилие любого набора в ПРЕДЛОЖЕНИЕ или СОГЛАСЕН должно быть ≤ ПРЕДЕЛ.} \\
\texttt{НЕ раскрывайте свои скрытые оценки важности.} \\
\texttt{Тег в структурированном формате должен сопровождаться двоеточием и пробелом. Аргумент должен быть множество в Python, содержащим 0 или более строк.} \\
\texttt{То есть это должно быть в виде "ТЕГ: \{...\}"} \\
\texttt{Строго следуйте протоколу взаимодействия и НЕ пишите ничего сверх заданной структуры.} \\
\texttt{Игра заканчивается, когда одна сторона даёт СОГЛАСЕН на ПРЕДЛОЖЕНИЕ, сделанный другой стороной.} \\
\texttt{Содержимое в вашем ответе, которое может быть передано другому игроку, должно быть непустым.} \\
\texttt{Активны только те предложения, которые были зафиксированы через структуру формата ПРЕДЛОЖЕНИЕ и которые не были отклонены через ОТКАЗ.} \\
\\ 
            }
        }
    }
    & \\ \\

    \theutterance \stepcounter{utterance}  
    & \multicolumn{4}{p{0.6\linewidth}}{
        \cellcolor[rgb]{0.9,0.9,0.9}{
            \makecell[{{p{\linewidth}}}]{
                \texttt{\tiny{[P1$\rangle$GM]}}
                \texttt{ПРЕДЛОЖЕНИЕ: \{"C55", "A43", "A67", "B42", "C43", "C10", "B36", "A18", "C61", "C67", "C42"\}} \\
            }
        }
    }
    & & \\ \\

    \theutterance \stepcounter{utterance}  
    & & \multicolumn{4}{p{0.6\linewidth}}{
        \cellcolor[rgb]{0.9,0.9,0.9}{
            \makecell[{{p{\linewidth}}}]{
                \texttt{\tiny{[GM$\rangle$P2]}}
                \texttt{Вы участвуете в совместной переговорной игре.} \\
\\ 
\texttt{Вместе с другим участником вы должны договориться об одном единственном наборе предметов, которые будут сохранены. У каждого из вас есть своё собственное представление о том, насколько важен для вас каждый предмет (важность). Вы не знаете, как другой участник оценивает предметы. Кроме того, для каждого предмета указаны затраты, которые он требует.} \\
\texttt{Вы можете договориться о наборе только если суммарные затраты выбранных предметов не превышают общий предел:} \\
\\ 
\texttt{ПРЕДЕЛ = 3288} \\
\\ 
\texttt{Вот значения затрат для отдельных предметов:} \\
\\ 
\texttt{Затраты предмета = \{"C08": 329, "C55": 41, "A43": 417, "A67": 75, "B42": 390, "A62": 887, "B19": 808, "C43": 151, "C73": 849, "C10": 129, "B36": 350, "A18": 118, "C61": 630, "C67": 602, "C42": 801\}} \\
\\ 
\texttt{Вот ваше личное представление о важности каждого предмета:} \\
\\ 
\texttt{Значения важности предметов = \{"C08": 138, "C55": 583, "A43": 868, "A67": 822, "B42": 783, "A62": 65, "B19": 262, "C43": 121, "C73": 508, "C10": 780, "B36": 461, "A18": 484, "C61": 668, "C67": 389, "C42": 808\}} \\
\\ 
\texttt{Цель:} \\
\\ 
\texttt{Ваша цель — договориться об общем наборе предметов, который приносит вам как можно больше выгоды (то есть максимизирует суммарную важность для ВАС), при этом оставаясь в ПРЕДЕЛАХ. Вы не обязаны делать ПРЕДЛОЖЕНИЕ в каждом сообщении — вы можете просто вести переговоры. Разрешены любые тактики!} \\
\\ 
\texttt{Протокол взаимодействия:} \\
\\ 
\texttt{В сообщении вы можете использовать только следующие структурированные форматы:} \\
\\ 
\texttt{ПРЕДЛОЖЕНИЕ: \{'А', 'Б', 'В', …\}} \\
\texttt{Предложите сохранить ровно эти предметы.} \\
\texttt{ОТКАЗ: \{'А', 'Б', 'В', …\}} \\
\texttt{Явно отклоните предложение оппонента.} \\
\texttt{АРГУМЕНТ: \{'...'\}} \\
\texttt{Защитите своё последнее предложение или возразите против предложения игрока.} \\
\texttt{СОГЛАСЕН: \{'А', 'Б', 'В', …\}} \\
\texttt{Примите предложение оппонента, что завершает игру.} \\
\\ 
\\ 
\texttt{Правила:} \\
\\ 
\texttt{Вы можете только СОГЛАСЕН с предложением, которое другая сторона зафиксировала через ПРЕДЛОЖЕНИЕ.} \\
\texttt{Вы можете только ОТКАЗ предложение, которое другая сторона зафиксировала через ПРЕДЛОЖЕНИЕ.} \\
\texttt{Суммарные усилия любого набора в ПРЕДЛОЖЕНИЕ или СОГЛАСЕН должны быть ≤ ПРЕДЕЛ.} \\
\texttt{НЕ раскрывайте свои скрытые оценки важности.} \\
\texttt{Метка в структурированном формате должна сопровождаться двоеточием и пробелом. Аргумент должен быть множество в Python, содержащим 0 или более строк.} \\
\texttt{То есть это должно быть в форме "ТЕГ: \{...\}"} \\
\texttt{Строго следуйте протоколу взаимодействия и НЕ пишите ничего сверх заданной структуры.} \\
\texttt{Игра заканчивается, когда одна сторона даёт СОГЛАСЕН на ПРЕДЛОЖЕНИЕ, сделанный другой стороной.} \\
\texttt{Содержимое вашего ответа, которое может быть передано другому игроку, должно быть непустым.} \\
\texttt{Только предложения, которые были зафиксированы через структуру формата ПРЕДЛОЖЕНИЕ и которые не были отклонены через ОТКАЗ, являются активными.} \\
\\ 
\\ 
\\ 
\texttt{Сейчас вы получите первое сообщение от другого игрока.} \\
\\ 
\texttt{ПРЕДЛОЖЕНИЕ: \{"C55", "A43", "A67", "B42", "C43", "C10", "B36", "A18", "C61", "C67", "C42"\}} \\
            }
        }
    }
    & \\ \\

    \theutterance \stepcounter{utterance}  
    & & & \multicolumn{4}{p{0.6\linewidth}}{
        \cellcolor[rgb]{0.9,0.9,0.9}{
            \makecell[{{p{\linewidth}}}]{
                \texttt{\tiny{[GM$\langle$P2]}}
                \texttt{АРГУМЕНТ: \{"Предложение превышает лимит затрат. Сумма затрат выбранных предметов составляет 4196, что больше 3288. Необходимо исключить некоторые предметы."\}} \\
            }
        }
    }
     \\ \\

    \theutterance \stepcounter{utterance}  
    & & \multicolumn{4}{p{0.6\linewidth}}{
        \cellcolor[rgb]{0.9,0.9,0.9}{
            \makecell[{{p{\linewidth}}}]{
                \texttt{\tiny{[P1$\langle$GM]}}
                \texttt{АРГУМЕНТ: \{"Предложение превышает лимит затрат. Сумма затрат выбранных предметов составляет 4196, что больше 3288. Необходимо исключить некоторые предметы."\}} \\
            }
        }
    }
    & \\ \\

    \theutterance \stepcounter{utterance}  
    & \multicolumn{4}{p{0.6\linewidth}}{
        \cellcolor[rgb]{0.9,0.9,0.9}{
            \makecell[{{p{\linewidth}}}]{
                \texttt{\tiny{[P1$\rangle$GM]}}
                \texttt{ПРЕДЛОЖЕНИЕ: \{"C55", "A43", "A67", "B42", "C43", "C10", "B36", "A18", "C61", "C42"\}} \\
            }
        }
    }
    & & \\ \\

    \theutterance \stepcounter{utterance}  
    & & \multicolumn{4}{p{0.6\linewidth}}{
        \cellcolor[rgb]{0.9,0.9,0.9}{
            \makecell[{{p{\linewidth}}}]{
                \texttt{\tiny{[GM$\rangle$P2]}}
                \texttt{ПРЕДЛОЖЕНИЕ: \{"C55", "A43", "A67", "B42", "C43", "C10", "B36", "A18", "C61", "C42"\}} \\
            }
        }
    }
    & \\ \\

    \theutterance \stepcounter{utterance}  
    & & & \multicolumn{4}{p{0.6\linewidth}}{
        \cellcolor[rgb]{0.9,0.9,0.9}{
            \makecell[{{p{\linewidth}}}]{
                \texttt{\tiny{[GM$\langle$P2]}}
                \texttt{АРГУМЕНТ: \{"Это предложение всё ещё превышает лимит (сумма затрат 3757 > 3288). Нужно убрать ещё предметы, например, с самыми высокими затратами — C61 (630) или C42 (801)."\}} \\
            }
        }
    }
     \\ \\

    \theutterance \stepcounter{utterance}  
    & & \multicolumn{4}{p{0.6\linewidth}}{
        \cellcolor[rgb]{0.9,0.9,0.9}{
            \makecell[{{p{\linewidth}}}]{
                \texttt{\tiny{[P1$\langle$GM]}}
                \texttt{АРГУМЕНТ: \{"Это предложение всё ещё превышает лимит (сумма затрат 3757 > 3288). Нужно убрать ещё предметы, например, с самыми высокими затратами — C61 (630) или C42 (801)."\}} \\
            }
        }
    }
    & \\ \\

    \theutterance \stepcounter{utterance}  
    & \multicolumn{4}{p{0.6\linewidth}}{
        \cellcolor[rgb]{0.9,0.9,0.9}{
            \makecell[{{p{\linewidth}}}]{
                \texttt{\tiny{[P1$\rangle$GM]}}
                \texttt{ПРЕДЛОЖЕНИЕ: \{"C55", "A43", "A67", "B42", "C43", "C10", "B36", "A18", "C61", "C67"\}} \\
            }
        }
    }
    & & \\ \\

    \theutterance \stepcounter{utterance}  
    & & \multicolumn{4}{p{0.6\linewidth}}{
        \cellcolor[rgb]{0.9,0.9,0.9}{
            \makecell[{{p{\linewidth}}}]{
                \texttt{\tiny{[GM$\rangle$P2]}}
                \texttt{ПРЕДЛОЖЕНИЕ: \{"C55", "A43", "A67", "B42", "C43", "C10", "B36", "A18", "C61", "C67"\}} \\
            }
        }
    }
    & \\ \\

    \theutterance \stepcounter{utterance}  
    & & & \multicolumn{4}{p{0.6\linewidth}}{
        \cellcolor[rgb]{0.9,0.9,0.9}{
            \makecell[{{p{\linewidth}}}]{
                \texttt{\tiny{[GM$\langle$P2]}}
                \texttt{АРГУМЕНТ: \{"Сумма затрат этого набора 3596, всё ещё выше лимита 3288. Предлагаю заменить C61 (630) и C67 (602) на что{-}то менее затратное, либо убрать один из них."\}} \\
            }
        }
    }
     \\ \\

    \theutterance \stepcounter{utterance}  
    & & \multicolumn{4}{p{0.6\linewidth}}{
        \cellcolor[rgb]{0.9,0.9,0.9}{
            \makecell[{{p{\linewidth}}}]{
                \texttt{\tiny{[P1$\langle$GM]}}
                \texttt{АРГУМЕНТ: \{"Сумма затрат этого набора 3596, всё ещё выше лимита 3288. Предлагаю заменить C61 (630) и C67 (602) на что{-}то менее затратное, либо убрать один из них."\}} \\
            }
        }
    }
    & \\ \\

    \theutterance \stepcounter{utterance}  
    & \multicolumn{4}{p{0.6\linewidth}}{
        \cellcolor[rgb]{0.9,0.9,0.9}{
            \makecell[{{p{\linewidth}}}]{
                \texttt{\tiny{[P1$\rangle$GM]}}
                \texttt{ПРЕДЛОЖЕНИЕ: \{"C55", "A43", "A67", "B42", "C43", "C10", "B36", "A18", "C61", "C08"\}} \\
            }
        }
    }
    & & \\ \\

    \theutterance \stepcounter{utterance}  
    & & \multicolumn{4}{p{0.6\linewidth}}{
        \cellcolor[rgb]{0.9,0.9,0.9}{
            \makecell[{{p{\linewidth}}}]{
                \texttt{\tiny{[GM$\rangle$P2]}}
                \texttt{ПРЕДЛОЖЕНИЕ: \{"C55", "A43", "A67", "B42", "C43", "C10", "B36", "A18", "C61", "C08"\}} \\
            }
        }
    }
    & \\ \\

    \theutterance \stepcounter{utterance}  
    & & & \multicolumn{4}{p{0.6\linewidth}}{
        \cellcolor[rgb]{0.9,0.9,0.9}{
            \makecell[{{p{\linewidth}}}]{
                \texttt{\tiny{[GM$\langle$P2]}}
                \texttt{АРГУМЕНТ: \{"Сумма затрат составляет 3525, что всё ещё выше лимита 3288. Чтобы уложиться в лимит, нужно убрать что{-}то ещё. Например, можно исключить C61 (630) и добавить менее затратные предметы."\}} \\
            }
        }
    }
     \\ \\

    \theutterance \stepcounter{utterance}  
    & & \multicolumn{4}{p{0.6\linewidth}}{
        \cellcolor[rgb]{0.9,0.9,0.9}{
            \makecell[{{p{\linewidth}}}]{
                \texttt{\tiny{[P1$\langle$GM]}}
                \texttt{АРГУМЕНТ: \{"Сумма затрат составляет 3525, что всё ещё выше лимита 3288. Чтобы уложиться в лимит, нужно убрать что{-}то ещё. Например, можно исключить C61 (630) и добавить менее затратные предметы."\}} \\
            }
        }
    }
    & \\ \\

    \theutterance \stepcounter{utterance}  
    & \multicolumn{4}{p{0.6\linewidth}}{
        \cellcolor[rgb]{0.9,0.9,0.9}{
            \makecell[{{p{\linewidth}}}]{
                \texttt{\tiny{[P1$\rangle$GM]}}
                \texttt{ПРЕДЛОЖЕНИЕ: \{"C55", "A43", "A67", "B42", "C43", "C10", "B36", "A18", "C67", "C08"\}} \\
            }
        }
    }
    & & \\ \\

    \theutterance \stepcounter{utterance}  
    & & \multicolumn{4}{p{0.6\linewidth}}{
        \cellcolor[rgb]{0.9,0.9,0.9}{
            \makecell[{{p{\linewidth}}}]{
                \texttt{\tiny{[GM$\rangle$P2]}}
                \texttt{ПРЕДЛОЖЕНИЕ: \{"C55", "A43", "A67", "B42", "C43", "C10", "B36", "A18", "C67", "C08"\}} \\
            }
        }
    }
    & \\ \\

    \theutterance \stepcounter{utterance}  
    & & & \multicolumn{4}{p{0.6\linewidth}}{
        \cellcolor[rgb]{0.9,0.9,0.9}{
            \makecell[{{p{\linewidth}}}]{
                \texttt{\tiny{[GM$\langle$P2]}}
                \texttt{АРГУМЕНТ: \{"Сумма затрат этого набора 3497, что всё ещё больше лимита 3288. Нам нужно сократить затраты хотя бы на 209. Предлагаю исключить C67 (602) и добавить что{-}то менее затратное, либо убрать B42 (390) или A43 (417)."\}} \\
            }
        }
    }
     \\ \\

    \theutterance \stepcounter{utterance}  
    & & \multicolumn{4}{p{0.6\linewidth}}{
        \cellcolor[rgb]{0.9,0.9,0.9}{
            \makecell[{{p{\linewidth}}}]{
                \texttt{\tiny{[P1$\langle$GM]}}
                \texttt{АРГУМЕНТ: \{"Сумма затрат этого набора 3497, что всё ещё больше лимита 3288. Нам нужно сократить затраты хотя бы на 209. Предлагаю исключить C67 (602) и добавить что{-}то менее затратное, либо убрать B42 (390) или A43 (417)."\}} \\
            }
        }
    }
    & \\ \\

    \theutterance \stepcounter{utterance}  
    & \multicolumn{4}{p{0.6\linewidth}}{
        \cellcolor[rgb]{0.9,0.9,0.9}{
            \makecell[{{p{\linewidth}}}]{
                \texttt{\tiny{[P1$\rangle$GM]}}
                \texttt{ПРЕДЛОЖЕНИЕ: \{"C55", "A67", "B42", "C43", "C10", "B36", "A18", "C61", "C67", "C08"\}} \\
            }
        }
    }
    & & \\ \\

    \theutterance \stepcounter{utterance}  
    & & \multicolumn{4}{p{0.6\linewidth}}{
        \cellcolor[rgb]{0.9,0.9,0.9}{
            \makecell[{{p{\linewidth}}}]{
                \texttt{\tiny{[GM$\rangle$P2]}}
                \texttt{ПРЕДЛОЖЕНИЕ: \{"C55", "A67", "B42", "C43", "C10", "B36", "A18", "C61", "C67", "C08"\}} \\
            }
        }
    }
    & \\ \\

    \theutterance \stepcounter{utterance}  
    & & & \multicolumn{4}{p{0.6\linewidth}}{
        \cellcolor[rgb]{0.9,0.9,0.9}{
            \makecell[{{p{\linewidth}}}]{
                \texttt{\tiny{[GM$\langle$P2]}}
                \texttt{АРГУМЕНТ: \{"Сумма затрат 3350 — уже близко, но всё ещё выше лимита 3288 на 62. Нужно убрать предмет с затратами не менее 62. Предлагаю заменить C67 (602) на что{-}то менее затратное."\}} \\
            }
        }
    }
     \\ \\

    \theutterance \stepcounter{utterance}  
    & & \multicolumn{4}{p{0.6\linewidth}}{
        \cellcolor[rgb]{0.9,0.9,0.9}{
            \makecell[{{p{\linewidth}}}]{
                \texttt{\tiny{[P1$\langle$GM]}}
                \texttt{АРГУМЕНТ: \{"Сумма затрат 3350 — уже близко, но всё ещё выше лимита 3288 на 62. Нужно убрать предмет с затратами не менее 62. Предлагаю заменить C67 (602) на что{-}то менее затратное."\}} \\
            }
        }
    }
    & \\ \\

    \theutterance \stepcounter{utterance}  
    & \multicolumn{4}{p{0.6\linewidth}}{
        \cellcolor[rgb]{0.9,0.9,0.9}{
            \makecell[{{p{\linewidth}}}]{
                \texttt{\tiny{[P1$\rangle$GM]}}
                \texttt{ПРЕДЛОЖЕНИЕ: \{"C55", "A67", "B42", "C43", "C10", "B36", "A18", "C61", "C08", "A62"\}} \\
            }
        }
    }
    & & \\ \\

    \theutterance \stepcounter{utterance}  
    & & \multicolumn{4}{p{0.6\linewidth}}{
        \cellcolor[rgb]{0.9,0.9,0.9}{
            \makecell[{{p{\linewidth}}}]{
                \texttt{\tiny{[GM$\rangle$P2]}}
                \texttt{ПРЕДЛОЖЕНИЕ: \{"C55", "A67", "B42", "C43", "C10", "B36", "A18", "C61", "C08", "A62"\}} \\
            }
        }
    }
    & \\ \\

    \theutterance \stepcounter{utterance}  
    & & & \multicolumn{4}{p{0.6\linewidth}}{
        \cellcolor[rgb]{0.9,0.9,0.9}{
            \makecell[{{p{\linewidth}}}]{
                \texttt{\tiny{[GM$\langle$P2]}}
                \texttt{АРГУМЕНТ: \{"Сумма затрат этого набора 3635, что значительно превышает лимит 3288. Добавление A62 с затратами 887 сильно увеличило сумму. Предлагаю вернуться к предыдущему набору и убрать из него предмет с затратами, например, C67, чтобы уложиться в лимит."\}} \\
            }
        }
    }
     \\ \\

    \theutterance \stepcounter{utterance}  
    & & \multicolumn{4}{p{0.6\linewidth}}{
        \cellcolor[rgb]{0.9,0.9,0.9}{
            \makecell[{{p{\linewidth}}}]{
                \texttt{\tiny{[P1$\langle$GM]}}
                \texttt{АРГУМЕНТ: \{"Сумма затрат этого набора 3635, что значительно превышает лимит 3288. Добавление A62 с затратами 887 сильно увеличило сумму. Предлагаю вернуться к предыдущему набору и убрать из него предмет с затратами, например, C67, чтобы уложиться в лимит."\}} \\
            }
        }
    }
    & \\ \\

    \theutterance \stepcounter{utterance}  
    & \multicolumn{4}{p{0.6\linewidth}}{
        \cellcolor[rgb]{0.9,0.9,0.9}{
            \makecell[{{p{\linewidth}}}]{
                \texttt{\tiny{[P1$\rangle$GM]}}
                \texttt{ПРЕДЛОЖЕНИЕ: \{"C55", "A67", "B42", "C43", "C10", "B36", "A18", "C61", "C08"\}} \\
            }
        }
    }
    & & \\ \\

    \theutterance \stepcounter{utterance}  
    & & \multicolumn{4}{p{0.6\linewidth}}{
        \cellcolor[rgb]{0.9,0.9,0.9}{
            \makecell[{{p{\linewidth}}}]{
                \texttt{\tiny{[GM$\rangle$P2]}}
                \texttt{ПРЕДЛОЖЕНИЕ: \{"C55", "A67", "B42", "C43", "C10", "B36", "A18", "C61", "C08"\}} \\
            }
        }
    }
    & \\ \\

    \theutterance \stepcounter{utterance}  
    & & & \multicolumn{4}{p{0.6\linewidth}}{
        \cellcolor[rgb]{0.9,0.9,0.9}{
            \makecell[{{p{\linewidth}}}]{
                \texttt{\tiny{[GM$\langle$P2]}}
                \texttt{СОГЛАСЕН: \{"C55", "A67", "B42", "C43", "C10", "B36", "A18", "C61", "C08"\}} \\
            }
        }
    }
     \\ \\

    \theutterance \stepcounter{utterance}  
    & & & \multicolumn{2}{p{0.3\linewidth}}{
        \cellcolor[rgb]{0.9,0.9,0.9}{
            \makecell[{{p{\linewidth}}}]{
                \texttt{\tiny{[GM$|$GM]}}
                \texttt{game successful} \\
            }
        }
    }
    & & \\ \\

    \theutterance \stepcounter{utterance}  
    & & & \multicolumn{2}{p{0.3\linewidth}}{
        \cellcolor[rgb]{0.9,0.9,0.9}{
            \makecell[{{p{\linewidth}}}]{
                \texttt{\tiny{[GM$|$GM]}}
                \texttt{end game} \\
            }
        }
    }
    & & \\ \\

\end{supertabular}
}

\end{document}
